\section{Approximate Join}\label{sec:technical-2:-approximate}\label{sec:approx-join}

The exact join drops almost every pair inside the kernel, but it still computes every pair.
A join of 10k probe rows with 100M corpus rows of 128 dimensions is $1.3\times10^{14}$ multiply-adds whatever the bound.
Relational engines face the same problem with large inputs, and they answer it with join algorithms that never compare most pairs.
Table~\ref{tab:join-analogy} lists the classic algorithms next to their vector counterparts.

\begin{table}[t]
    \centering\footnotesize
    \setlength{\tabcolsep}{3pt}
    \begin{tabular}{@{}p{0.17\columnwidth}p{0.43\columnwidth}p{0.33\columnwidth}@{}}
        \toprule
        \raggedright relational join & \raggedright vector counterpart & \raggedright examples \tabularnewline
        \midrule
        \raggedright block nested loop & \raggedright compare every probe tile with every corpus tile & \raggedright Sirius exact join (\S\ref{sec:exact-join}), Databricks \texttt{NEAREST BY}~\cite{DatabricksNearestByJoinBlog26}, DuckDB \texttt{vss\_join}~\cite{DuckDBVSS} \tabularnewline
        \raggedright index nested loop & \raggedright one search of a prebuilt index per probe row & \raggedright FAISS~\cite{JohnsonFaissGPU21}, cuVS~\cite{NvidiaCuVS26}, DuckDB HNSW join~\cite{DuckDBVSSBlog24}, MaxVec~\cite{MaxVec26} \tabularnewline
        \raggedright hash join & \raggedright partition the corpus at query time, route each probe row to its $n_p$ nearest partitions & \raggedright Sirius approximate join (this section), SemDeDup~\cite{AbbasSemDeDup23} \tabularnewline
        \bottomrule
    \end{tabular}
    \caption{Relational join algorithms and their vector counterparts.}
    \label{tab:join-analogy}
\end{table}

A vector-search library answers a join as an index nested-loop join, which pays off, as for relational data, when the index already exists and the probe side is small.
Our approximate join is a hash join whose hash function is k-means: it partitions the corpus at query time and sends each probe row to its nearest partitions (Figure~\ref{fig:hash-vs-approx}).

\input{hash-vs-approx}

\subsection{Build: Partitioning the Corpus}\label{subsec:approx-build}

\yxy{we use the word corpus a lot. Isn't this term more specific to vector search? Is it better to call it build table vs. probe table instead?}

Sirius partitions the corpus inside the statement that needs the partitions, as a hash join builds its table.
It samples 256 rows per centroid, or fewer when that sample does not fit beside what the GPU already holds, and fits cuVS's balanced k-means, the one its IVF indexes use~\cite{NvidiaCuVS26}.
The k-means is balanced because a partition is the unit of work below, and a few oversized partitions would leave most of the GPU idle.
Sirius then assigns every corpus row to its nearest centroid.
This assignment is a per-row top-1 join of the corpus with the centroids, which Sirius runs with the bounded kernel of Section~\ref{subsec:exact-kernel}.
Finally it writes the corpus in partition order, one contiguous run per partition, coded as in Section~\ref{subsec:exact-codes}, with a map back to the row ids.
The partitions stay in the session for later statements.

The number of partitions $N$ trades build time for join time.
Sirius starts from $N_{\mathrm{ref}}=\sqrt{n}$, at most 1\,024, and considers its doublings up to $16N_{\mathrm{ref}}$ that leave at least 256 rows per partition.
It picks the count with the smallest estimated build plus join cost for the estimated number of probe rows (Section~\ref{subsec:integration-plan}).
A plain-SQL join is exact by default; a user asks for an approximate answer with \texttt{SET vector\_join\_probes = }$n_p$, as with pgvector's \texttt{ivfflat.probes}.
Sirius scales $n_p$ by $\sqrt{N/N_{\mathrm{ref}}}$, because recall follows $n_p/\sqrt{N}$ empirically: on Deep10M, one setting keeps recall@10 between 0.956 and 0.964 for every count from 256 to 16\,384.

\subsection{Probe: Routing Each Probe Row\yxy{maybe more like co-partitioning the probe table?}}\label{subsec:approx-probe}

Routing is again a join: probe rows are multiplied with the $N$ centroids in one GEMM\yxy{it's actually a smaller vector join. We chose to implement it using GEMM but that's one implementation decision, not the only one. should clarify}, and a per-row top-$n_p$ keeps each row's $n_p$ nearest partitions.
Sorting the resulting (probe row, partition) pairs by partition gives one \emph{slice} per partition: the partition and every probe row routed to it.
Here the vector join departs from an equi-join.
A probe row visits $n_p$ buckets instead of one, because its nearest neighbors can lie just across a partition boundary.
A row of the answer is missed only when its partition is not among those its probe row visits.
So $n_p$ trades recall for time, and $n_p=N$ returns the exact answer.

\subsection{Join: Each Partition Once}\label{subsec:approx-join}

\yxy{my current understanding is: after co-partitioning, we simply divide the original vector join problem into 1024 (or whatever the number is) smaller vector join problems that hopefully fits in GPU memory.
We can use any vector join algorithm to tackle this small vector join. In our current implementation, we decide to use our exact vector join algorithm to implement it. But it's possible to improve this in future work.}

Each slice is a small exact join, run with the kernel of Section~\ref{sec:exact-join}.
A partition is read once per batch of probe rows and multiplied in one GEMM with every probe row of the batch routed to it, whereas an index nested-loop join reads it once for every probe row that visits it; with 10k probe rows, 256 partitions and $n_p=8$, each read serves about 300 probe rows.
FAISS's \texttt{big\_batch\_search} processes an IVF search list by list in the same order~\cite{DouzeFaiss24}; Sirius makes this order the join's plan.

\yxy{is below the same optimization as already discussed in exact join section? If so, no need to repeat since we are just reusing our exact join for per-partition join.}
Routing also supplies the bound.
A first sweep joins each probe row with its nearest partition, and the $k$-th distance found there, re-scored in fp32, seeds the row's bound (Section~\ref{subsec:exact-bound}).
A second sweep runs the bounded kernel over every routed slice and lowers each row's bound as its slices finish.
Slices too small to pay for a kernel launch of their own are launched together.
The codes and the bound lose nothing, so the answer equals the exact join restricted to the probed partitions: routing is the only approximation, and every reduction of Section~\ref{subsec:reduction-modes} carries over.
